\section{How to write long-form proofs in LaTeX}\label{proofs: sec: latex}\hypertarget{proofs.latex}{}

\doHeader{Long-form proofs in LaTeX}

Instructions for getting started with LaTeX are in \bookLink{latex.homeworks}{Appendix B of the book}.
This section assumes you've read them.
This one is about how to use LaTeX (and the macros in the homework templates) to write long-form proofs.

% As usual, your LaTeX document should start with the standard header,
% which includes the macros file for the course, similar to what's shown in Figure~\ref{proofs: fig: header}.
% You can find macros file on Blackboard.
% It provides some LaTeX macros and environments for writing long-form proofs.
% We describe these next.

% \begin{figure}
% \begin{framed}\small
% % \vspace*{-10pt}
% \begin{verbatim}
% \documentclass[11pt]{article}

% ../macros/macros.tex

% \begin{document}

% ...

% \end{document}
% \end{verbatim}
% \vspace*{-10pt}
% \end{framed}
% \caption{The header and closing statement of a LaTeX document.
%   Beware that the macro file may be named something other than ``\texttt{macros.tex}'',
%   in which case you'll have to change the argument to the \texttt{\textbackslash input} statement accordingly.
%   For the homework, the LaTeX templates will already have this part,
%   so if you are using them you don't have to worry about this part.
% }\label{proofs: fig: header}
% \end{figure}

Within the body of the document, you can state a theorem or lemma using the \texttt{theorem} or
\texttt{lemma} environments. % , as shown in Figure~\ref{proofs: fig: theorem}.
LaTeX will automatically number the theorem or lemma for you.
% \begin{figure}
% \begin{framed}\small
% % \vspace*{-10pt}
% \begin{verbatim}
% \begin{theorem}\label{thm: primes}
%   There are infinitely many primes.
% \end{theorem}
% \end{verbatim}
% \vspace*{-10pt}
% \end{framed}
% \caption{Use the \texttt{theorem} environment to state a theorem.
%   LaTeX will automatically number the theorem.
%   You can record the theorem number using the \texttt{\textbackslash label} command,
%   then use the number later via the \texttt{\textbackslash ref} command.
% }\label{proofs: fig: theorem}
% \end{figure}
Follow your theorem or lemma statement by a long-form proof, using the \texttt{longFormProof} environment.
With that, start each step with a \texttt{\textbackslash step} command. % as shown in Figure~\ref{proofs: fig: longFormProof}.
% %
% \begin{figure}
% \begin{framed}\small
% % \vspace*{-10pt}
% \begin{verbatim}
% \begin{longFormProof}
%   \step Suppose for contradiction that there are finitely many primes.
%   \step Let $p_1, p_2, \ldots, p_n$ be a list of all the primes (larger than 1).
%   \step ...
% \end{longFormProof}
% \end{verbatim}
% \vspace*{-10pt}
% \end{framed}
% \caption{Use the \texttt{longFormProof} environment and the \texttt{\textbackslash step} command
%   to make a long-form proof.}\label{proofs: fig: longFormProof}
% \end{figure}
%
To introduce an inner block, use the \texttt{block} environment. % as shown in Figure~\ref{proofs: fig: blocks}.
Figure~\ref{proofs: fig: blocks} illustrates these environments and the \texttt{\textbackslash step} command.
%
\begin{figure}
\begin{framed}\small
% \vspace*{-10pt}
\begin{verbatim}
\begin{theorem}\label{thm: bad cases}
  For every real number, its square is positive. 
  That is, $\forall x\in\R,~x^2 > 0$.
\end{theorem}

\begin{longFormProof}
  \step Let $x\in \R$ be an arbitrary real number.

  \begin{block}
    \step \underline{Case 1.} Suppose $x$ is positive ($x > 0$).
    \step Then $x^2$ is a product of two positive numbers, so $x^2$ is positive.
  \end{block}

  \begin{block}
    \step \underline{Case 2.} Suppose $x$ is negative ($x < 0$).
    \step Then $x^2$ is a product of two negative numbers, so $x^2$ is positive.
  \end{block}

  \step In either Case 1 or 2, $x^2$ is positive, so $x^2$ is positive.
\end{longFormProof}
\end{verbatim}
\vspace*{-10pt}
\end{framed}
\caption{How to use the \texttt{theorem}, \texttt{longFormProof}, and \texttt{block} environments,
  and the \texttt{\textbackslash step} command.
  See Exercise~\ref{proofs: exercise: cases} for the result.}\label{proofs: fig: blocks}
\end{figure}
%

Optionally, you can use a \texttt{\textbackslash label} command to 
record the theorem or lemma number, using a name of your choice.
Then you can recall the number by using the \texttt{\textbackslash ref} command.
Specifically, \texttt{\textbackslash ref\{thm:~bad cases\}} will reproduce the number
(assuming you use the same name as you gave to \texttt{\textbackslash label}).
You can use \texttt{\textbackslash label} and \texttt{\textbackslash ref}
to label and refer to any type of numbered object in your document,
such as a proof step or block.
When the numbers change, LaTeX automatically updates the references.
(It may take two iterations of running LaTeX for the references to fully update.)

To introduce a lemma inside a long form proof,
use the \texttt{lemma} environment to state the lemma,
then use the \texttt{longFormSubProof} environment for the proof of the lemma.
An example is shown in Figure~\ref{proofs: fig: lemma}.

\begin{figure}
\begin{framed}\small
\begin{verbatim}
\begin{theorem}
  Any finite, connected graph $G$ that has Property $A$ also has Property $B$.
\end{theorem}

\begin{longFormProof}[First long-form proof.]
  \step Let $G$ be an arbitrary finite, connected graph having Property $A$.

  \step Before we give the proof, we prove the following utility lemma:

  \begin{lemma}\label{lemma: high low}
    Let $v$ be any vertex in $G$ that has a neighbor $x$ with larger weight
    ($w(x) > w(v)$). Then $v$ has a neighbor $y$ with a smaller weight ($w(y) < w(v)$).
  \end{lemma}

  \begin{longFormSubProof}[Proof of lemma.]

    \step ...

  \end{longFormSubProof}

  \begin{block} \label{A}
    \step Now assume for contradiction that $G$ does not have Property $B$.
    
    \step ...

    \step So, by Lemma~\ref{lemma: high low} ...

    \step This is a contradiction (as $G$ is finite).  
  \end{block}

  \step By Block~\ref{A}, $G$ must have Property $B$.

\end{longFormProof}
\end{verbatim}
\vspace*{-10pt}
\end{framed}
\caption{Making a lemma with the \texttt{lemma} and \texttt{longFormSubProof} environments.
  See the solution to Problem~\ref{proofs: prob: with lemma} in Section~\ref{proofs: sec: exercises} for the output.
}\label{proofs: fig: lemma}
\end{figure}

%%% Local Variables:
%%% mode: latex
%%% TeX-master: "long_form_proofs"
%%% End:
